\editionPage[\pdfbookmark[0]{Défis d’approfondissement}{bm-defis-debut-parent}]{defis-debut}{Vingt-et-un défis}{Un défi par chapitre · Parcours facultatif}{Chercher avant de consulter les deux indications.}
\begin{blocEdition}{defis-debut}
Ces problèmes demandent davantage d’initiative. Ils utilisent les outils du livre ; les prolongements sont définis dans l’énoncé. Les défis sont facultatifs et approfondissent les notions du programme.\par\bigskip
\begin{savoirfaire}[Une recherche qui avance]Faire un dessin, tester un petit cas, organiser des essais dans un tableau, garder ses traces. Ouvrir la première indication seulement après une vraie recherche ; revenir au problème avant de lire la seconde.\end{savoirfaire}
\begin{savoirfaire}[Après la solution]Repérer l’idée décisive et la refaire sur une feuille vierge, puis raconter la solution à quelqu’un. Expliquer pourquoi cela fonctionne compte autant que trouver.\end{savoirfaire}
\end{blocEdition}
\editionPage{defi-ex-1}{D01 · Défi}{Compter jusqu’au milliard}{Approfondissement · 20 à 45 minutes indicatives}
\begin{blocEdition}{defi-ex-1}
On compte à voix haute, un nombre par seconde, sans jamais s’arrêter. Combien de temps faut-il pour compter jusqu’à un million ? Répondre en jours, heures, minutes et secondes. Estimer ensuite, en années, le temps nécessaire pour compter jusqu’à un milliard (la calculatrice est autorisée pour cette dernière question).\par\vspace{5mm}{\otLight\fontsize{8}{10}\selectfont\color{otMuted}Espace de recherche}\par\par\vspace{7mm}{\color{otRule!65}\hrule height .25pt}\par\vspace{7mm}{\color{otRule!65}\hrule height .25pt}\par\vspace{7mm}{\color{otRule!65}\hrule height .25pt}\par\vspace{7mm}{\color{otRule!65}\hrule height .25pt}\par\vspace{7mm}{\color{otRule!65}\hrule height .25pt}\par\vspace{7mm}{\color{otRule!65}\hrule height .25pt}
\end{blocEdition}
\editionPage{defi-ex-2}{D02 · Défi}{Entre 3,7 et 3,8}{Approfondissement · 20 à 45 minutes indicatives}
\begin{blocEdition}{defi-ex-2}
Combien y a-t-il de nombres qui s’écrivent avec exactement deux décimales et sont strictement compris entre $3{,}7$ et $3{,}8$ ? Combien avec au plus trois décimales ? Peut-on écrire la liste complète de tous les nombres décimaux compris entre $3{,}7$ et $3{,}8$ ?\par\vspace{5mm}{\otLight\fontsize{8}{10}\selectfont\color{otMuted}Espace de recherche}\par\par\vspace{7mm}{\color{otRule!65}\hrule height .25pt}\par\vspace{7mm}{\color{otRule!65}\hrule height .25pt}\par\vspace{7mm}{\color{otRule!65}\hrule height .25pt}\par\vspace{7mm}{\color{otRule!65}\hrule height .25pt}\par\vspace{7mm}{\color{otRule!65}\hrule height .25pt}\par\vspace{7mm}{\color{otRule!65}\hrule height .25pt}
\end{blocEdition}
\editionPage{defi-ex-3}{D03 · Défi}{Somme 5, produit 6,09}{Approfondissement · 20 à 45 minutes indicatives}
\begin{blocEdition}{defi-ex-3}
Trouver deux nombres décimaux, ayant chacun un seul chiffre après la virgule, dont la somme vaut 5 et le produit vaut $6{,}09$.\par\vspace{5mm}{\otLight\fontsize{8}{10}\selectfont\color{otMuted}Espace de recherche}\par\par\vspace{7mm}{\color{otRule!65}\hrule height .25pt}\par\vspace{7mm}{\color{otRule!65}\hrule height .25pt}\par\vspace{7mm}{\color{otRule!65}\hrule height .25pt}\par\vspace{7mm}{\color{otRule!65}\hrule height .25pt}\par\vspace{7mm}{\color{otRule!65}\hrule height .25pt}\par\vspace{7mm}{\color{otRule!65}\hrule height .25pt}
\end{blocEdition}
\editionPage{defi-ex-4}{D04 · Défi}{Les restes en escalier}{Approfondissement · 20 à 45 minutes indicatives}
\begin{blocEdition}{defi-ex-4}
Trouver le plus petit nombre entier qui, divisé par 2, 3, 4, 5 et 6, donne chaque fois un reste égal au diviseur moins 1 (reste 1 par 2, reste 2 par 3, etc.). Trouver ensuite le plus petit nombre qui vérifie ces conditions et qui est, en plus, divisible par 7.\par\vspace{5mm}{\otLight\fontsize{8}{10}\selectfont\color{otMuted}Espace de recherche}\par\par\vspace{7mm}{\color{otRule!65}\hrule height .25pt}\par\vspace{7mm}{\color{otRule!65}\hrule height .25pt}\par\vspace{7mm}{\color{otRule!65}\hrule height .25pt}\par\vspace{7mm}{\color{otRule!65}\hrule height .25pt}\par\vspace{7mm}{\color{otRule!65}\hrule height .25pt}\par\vspace{7mm}{\color{otRule!65}\hrule height .25pt}
\end{blocEdition}
\editionPage{defi-ex-5}{D05 · Défi}{Les pains égyptiens}{Approfondissement · 20 à 45 minutes indicatives}
\begin{blocEdition}{defi-ex-5}
Il y a près de quatre mille ans, les scribes égyptiens partageaient des pains en n’utilisant que des parts de la forme un demi, un quart, un huitième… Comment partager 5 pains identiques entre 8 personnes, en coupant le moins possible, pour que chacun reçoive $\frac58$ de pain ? Et 7 pains entre 8 personnes ?\par\vspace{5mm}{\otLight\fontsize{8}{10}\selectfont\color{otMuted}Espace de recherche}\par\par\vspace{7mm}{\color{otRule!65}\hrule height .25pt}\par\vspace{7mm}{\color{otRule!65}\hrule height .25pt}\par\vspace{7mm}{\color{otRule!65}\hrule height .25pt}\par\vspace{7mm}{\color{otRule!65}\hrule height .25pt}\par\vspace{7mm}{\color{otRule!65}\hrule height .25pt}\par\vspace{7mm}{\color{otRule!65}\hrule height .25pt}
\end{blocEdition}
\editionPage{defi-ex-6}{D06 · Défi}{Une somme qui se simplifie}{Approfondissement · 20 à 45 minutes indicatives}
\begin{blocEdition}{defi-ex-6}
Calculer $\frac12-\frac13$, puis $\frac13-\frac14$, puis $\frac14-\frac15$. En déduire, sans dénominateur commun compliqué, la somme $\frac12+\frac16+\frac1{12}+\frac1{20}$.\par\vspace{5mm}{\otLight\fontsize{8}{10}\selectfont\color{otMuted}Espace de recherche}\par\par\vspace{7mm}{\color{otRule!65}\hrule height .25pt}\par\vspace{7mm}{\color{otRule!65}\hrule height .25pt}\par\vspace{7mm}{\color{otRule!65}\hrule height .25pt}\par\vspace{7mm}{\color{otRule!65}\hrule height .25pt}\par\vspace{7mm}{\color{otRule!65}\hrule height .25pt}\par\vspace{7mm}{\color{otRule!65}\hrule height .25pt}
\end{blocEdition}
\editionPage{defi-ex-7}{D07 · Défi}{Hausse puis baisse}{Approfondissement · 20 à 45 minutes indicatives}
\begin{blocEdition}{defi-ex-7}
Un vélo coûte 100 €. Son prix augmente de 10~\%, puis le nouveau prix baisse de 10~\%. Revient-on à 100 € ? Que se passe-t-il si la baisse a lieu d’abord ? Et avec une hausse de 20~\% suivie d’une baisse de 20~\% ?\par\vspace{5mm}{\otLight\fontsize{8}{10}\selectfont\color{otMuted}Espace de recherche}\par\par\vspace{7mm}{\color{otRule!65}\hrule height .25pt}\par\vspace{7mm}{\color{otRule!65}\hrule height .25pt}\par\vspace{7mm}{\color{otRule!65}\hrule height .25pt}\par\vspace{7mm}{\color{otRule!65}\hrule height .25pt}\par\vspace{7mm}{\color{otRule!65}\hrule height .25pt}\par\vspace{7mm}{\color{otRule!65}\hrule height .25pt}
\end{blocEdition}
\editionPage{defi-ex-8}{D08 · Défi}{Les tables du banquet}{Approfondissement · 20 à 45 minutes indicatives}
\begin{blocEdition}{defi-ex-8}
Pour un banquet, on aligne des tables carrées : une table seule accueille 4 personnes, deux tables accolées en accueillent 6, trois tables en accueillent 8. Combien de personnes peut-on asseoir avec 10 tables ? Combien de tables faut-il pour 50 invités ?\par\vspace{5mm}{\otLight\fontsize{8}{10}\selectfont\color{otMuted}Espace de recherche}\par\par\vspace{7mm}{\color{otRule!65}\hrule height .25pt}\par\vspace{7mm}{\color{otRule!65}\hrule height .25pt}\par\vspace{7mm}{\color{otRule!65}\hrule height .25pt}\par\vspace{7mm}{\color{otRule!65}\hrule height .25pt}\par\vspace{7mm}{\color{otRule!65}\hrule height .25pt}\par\vspace{7mm}{\color{otRule!65}\hrule height .25pt}
\end{blocEdition}
\editionPage{defi-ex-9}{D09 · Défi}{Les engrenages}{Approfondissement · 20 à 45 minutes indicatives}
\begin{blocEdition}{defi-ex-9}
Une petite roue dentée de 12 dents entraîne une grande roue de 36 dents. Quand la petite roue fait 15 tours, combien de tours fait la grande ? La grande roue entraîne à son tour une roue de 24 dents : combien de tours fait cette troisième roue ?\par\vspace{5mm}{\otLight\fontsize{8}{10}\selectfont\color{otMuted}Espace de recherche}\par\par\vspace{7mm}{\color{otRule!65}\hrule height .25pt}\par\vspace{7mm}{\color{otRule!65}\hrule height .25pt}\par\vspace{7mm}{\color{otRule!65}\hrule height .25pt}\par\vspace{7mm}{\color{otRule!65}\hrule height .25pt}\par\vspace{7mm}{\color{otRule!65}\hrule height .25pt}\par\vspace{7mm}{\color{otRule!65}\hrule height .25pt}
\end{blocEdition}
\editionPage{defi-ex-10}{D10 · Défi}{Une ficelle autour de la Terre}{Approfondissement · 20 à 45 minutes indicatives}
\begin{blocEdition}{defi-ex-10}
On entoure l’équateur, long d’environ $40\,000$ km, avec une ficelle bien tendue. On allonge la ficelle de 1 m et on la soulève régulièrement tout autour. Un chat peut-il passer dessous ? Faire le même calcul avec un ballon de 70 cm de tour.\par\vspace{5mm}{\otLight\fontsize{8}{10}\selectfont\color{otMuted}Espace de recherche}\par\par\vspace{7mm}{\color{otRule!65}\hrule height .25pt}\par\vspace{7mm}{\color{otRule!65}\hrule height .25pt}\par\vspace{7mm}{\color{otRule!65}\hrule height .25pt}\par\vspace{7mm}{\color{otRule!65}\hrule height .25pt}\par\vspace{7mm}{\color{otRule!65}\hrule height .25pt}\par\vspace{7mm}{\color{otRule!65}\hrule height .25pt}
\end{blocEdition}
\editionPage{defi-ex-11}{D11 · Défi}{Le carré penché}{Approfondissement · 20 à 45 minutes indicatives}
\begin{blocEdition}{defi-ex-11}
Sur un quadrillage de carreaux de 1 cm, on trace un carré penché dont les sommets sont des nœuds du quadrillage. Il est inscrit dans un carré de 4 cm de côté. Calculer son aire, sans formule nouvelle, en découpant et en recomposant.\label{fig-defi-11}\figureGeo{defi-carre}
\par\vspace{5mm}{\otLight\fontsize{8}{10}\selectfont\color{otMuted}Espace de recherche}\par\par\vspace{7mm}{\color{otRule!65}\hrule height .25pt}\par\vspace{7mm}{\color{otRule!65}\hrule height .25pt}\par\vspace{7mm}{\color{otRule!65}\hrule height .25pt}\par\vspace{7mm}{\color{otRule!65}\hrule height .25pt}\par\vspace{7mm}{\color{otRule!65}\hrule height .25pt}\par\vspace{7mm}{\color{otRule!65}\hrule height .25pt}
\end{blocEdition}
\editionPage{defi-ex-12}{D12 · Défi}{Les deux feux}{Approfondissement · 20 à 45 minutes indicatives}
\begin{blocEdition}{defi-ex-12}
Un feu clignote toutes les 45 secondes et un autre toutes les 1 min 20 s. Ils clignotent ensemble à 8 h 00 min 00 s. À quelle heure clignoteront-ils de nouveau ensemble ?\par\vspace{5mm}{\otLight\fontsize{8}{10}\selectfont\color{otMuted}Espace de recherche}\par\par\vspace{7mm}{\color{otRule!65}\hrule height .25pt}\par\vspace{7mm}{\color{otRule!65}\hrule height .25pt}\par\vspace{7mm}{\color{otRule!65}\hrule height .25pt}\par\vspace{7mm}{\color{otRule!65}\hrule height .25pt}\par\vspace{7mm}{\color{otRule!65}\hrule height .25pt}\par\vspace{7mm}{\color{otRule!65}\hrule height .25pt}
\end{blocEdition}
\editionPage{defi-ex-13}{D13 · Défi}{Combien de points ?}{Approfondissement · 20 à 45 minutes indicatives}
\begin{blocEdition}{defi-ex-13}
Les points A et B sont tels que $AB=6$ cm. On cherche les points situés à 4 cm de A et à $r$ cm de B. Combien y a-t-il de tels points quand $r$ vaut 1, 2, 5, 10, puis 11 ?\label{fig-defi-13}\figureGeo{defi-cercles}
\par\vspace{5mm}{\otLight\fontsize{8}{10}\selectfont\color{otMuted}Espace de recherche}\par\par\vspace{7mm}{\color{otRule!65}\hrule height .25pt}\par\vspace{7mm}{\color{otRule!65}\hrule height .25pt}\par\vspace{7mm}{\color{otRule!65}\hrule height .25pt}\par\vspace{7mm}{\color{otRule!65}\hrule height .25pt}\par\vspace{7mm}{\color{otRule!65}\hrule height .25pt}\par\vspace{7mm}{\color{otRule!65}\hrule height .25pt}
\end{blocEdition}
\editionPage{defi-ex-14}{D14 · Défi}{Le centre perdu}{Approfondissement · 20 à 45 minutes indicatives}
\begin{blocEdition}{defi-ex-14}
Une assiette ronde laisse sur une feuille un cercle, mais son centre n’est pas marqué. Comment retrouver ce centre avec une règle et un compas ? Expliquer pourquoi la méthode fonctionne.\label{fig-defi-14}\figureGeo{defi-centre}
\par\vspace{5mm}{\otLight\fontsize{8}{10}\selectfont\color{otMuted}Espace de recherche}\par\par\vspace{7mm}{\color{otRule!65}\hrule height .25pt}\par\vspace{7mm}{\color{otRule!65}\hrule height .25pt}\par\vspace{7mm}{\color{otRule!65}\hrule height .25pt}\par\vspace{7mm}{\color{otRule!65}\hrule height .25pt}\par\vspace{7mm}{\color{otRule!65}\hrule height .25pt}\par\vspace{7mm}{\color{otRule!65}\hrule height .25pt}
\end{blocEdition}
\editionPage{defi-ex-15}{D15 · Défi}{Les aiguilles en mouvement}{Approfondissement · 20 à 45 minutes indicatives}
\begin{blocEdition}{defi-ex-15}
En une heure, la grande aiguille fait un tour complet et la petite aiguille avance d’une heure, soit $30^\circ$. Quel angle forment les deux aiguilles à 2 h 30 ? À 9 h 45 ?\par\vspace{5mm}{\otLight\fontsize{8}{10}\selectfont\color{otMuted}Espace de recherche}\par\par\vspace{7mm}{\color{otRule!65}\hrule height .25pt}\par\vspace{7mm}{\color{otRule!65}\hrule height .25pt}\par\vspace{7mm}{\color{otRule!65}\hrule height .25pt}\par\vspace{7mm}{\color{otRule!65}\hrule height .25pt}\par\vspace{7mm}{\color{otRule!65}\hrule height .25pt}\par\vspace{7mm}{\color{otRule!65}\hrule height .25pt}
\end{blocEdition}
\editionPage{defi-ex-16}{D16 · Défi}{Des triangles d’allumettes}{Approfondissement · 20 à 45 minutes indicatives}
\begin{blocEdition}{defi-ex-16}
On dispose de 12 allumettes identiques et on veut former un triangle en les utilisant toutes, sans les casser. Combien de triangles différents peut-on obtenir ? Lesquels sont isocèles ? équilatéraux ?\par\vspace{5mm}{\otLight\fontsize{8}{10}\selectfont\color{otMuted}Espace de recherche}\par\par\vspace{7mm}{\color{otRule!65}\hrule height .25pt}\par\vspace{7mm}{\color{otRule!65}\hrule height .25pt}\par\vspace{7mm}{\color{otRule!65}\hrule height .25pt}\par\vspace{7mm}{\color{otRule!65}\hrule height .25pt}\par\vspace{7mm}{\color{otRule!65}\hrule height .25pt}\par\vspace{7mm}{\color{otRule!65}\hrule height .25pt}
\end{blocEdition}
\editionPage{defi-ex-17}{D17 · Défi}{Deux miroirs parallèles}{Approfondissement · 20 à 45 minutes indicatives}
\begin{blocEdition}{defi-ex-17}
Les droites $(d_1)$ et $(d_2)$ sont parallèles et distantes de 3 cm. On construit le symétrique $M_1$ d’un point M par rapport à $(d_1)$, puis le symétrique $M_2$ de $M_1$ par rapport à $(d_2)$. Mesurer $MM_2$. Recommencer avec un autre point M. Que remarque-t-on ?\label{fig-defi-17}\figureGeo{defi-miroirs}
\par\vspace{5mm}{\otLight\fontsize{8}{10}\selectfont\color{otMuted}Espace de recherche}\par\par\vspace{7mm}{\color{otRule!65}\hrule height .25pt}\par\vspace{7mm}{\color{otRule!65}\hrule height .25pt}\par\vspace{7mm}{\color{otRule!65}\hrule height .25pt}\par\vspace{7mm}{\color{otRule!65}\hrule height .25pt}\par\vspace{7mm}{\color{otRule!65}\hrule height .25pt}\par\vspace{7mm}{\color{otRule!65}\hrule height .25pt}
\end{blocEdition}
\editionPage{defi-ex-18}{D18 · Défi}{La pyramide de cubes}{Approfondissement · 20 à 45 minutes indicatives}
\begin{blocEdition}{defi-ex-18}
On construit une pyramide pleine en empilant des étages carrés de cubes de 1 cm d’arête : 1 cube en haut, puis 2 sur 2, puis 3 sur 3, et ainsi de suite. Quel est le volume d’une pyramide de 4 étages ? De 5 étages ? Combien d’étages faut-il pour dépasser 200 cm$^3$ ?\label{fig-defi-18}\figureGeo{defi-pyramide}
\par\vspace{5mm}{\otLight\fontsize{8}{10}\selectfont\color{otMuted}Espace de recherche}\par\par\vspace{7mm}{\color{otRule!65}\hrule height .25pt}\par\vspace{7mm}{\color{otRule!65}\hrule height .25pt}\par\vspace{7mm}{\color{otRule!65}\hrule height .25pt}\par\vspace{7mm}{\color{otRule!65}\hrule height .25pt}\par\vspace{7mm}{\color{otRule!65}\hrule height .25pt}\par\vspace{7mm}{\color{otRule!65}\hrule height .25pt}
\end{blocEdition}
\editionPage{defi-ex-19}{D19 · Défi}{Le tableau à reconstituer}{Approfondissement · 20 à 45 minutes indicatives}
\begin{blocEdition}{defi-ex-19}
Dans un collège, 50 élèves de Sixième ont répondu à une enquête. On sait que 28 sont des filles, que 30 élèves déjeunent à la cantine et que 12 garçons déjeunent à la cantine. Reconstituer le tableau à double entrée (filles ou garçons ; cantine ou non). Quelle proportion des filles déjeune à la cantine ?\par\vspace{5mm}{\otLight\fontsize{8}{10}\selectfont\color{otMuted}Espace de recherche}\par\par\vspace{7mm}{\color{otRule!65}\hrule height .25pt}\par\vspace{7mm}{\color{otRule!65}\hrule height .25pt}\par\vspace{7mm}{\color{otRule!65}\hrule height .25pt}\par\vspace{7mm}{\color{otRule!65}\hrule height .25pt}\par\vspace{7mm}{\color{otRule!65}\hrule height .25pt}\par\vspace{7mm}{\color{otRule!65}\hrule height .25pt}
\end{blocEdition}
\editionPage{defi-ex-20}{D20 · Défi}{Les chaussettes dans le noir}{Approfondissement · 20 à 45 minutes indicatives}
\begin{blocEdition}{defi-ex-20}
Un tiroir contient 5 chaussettes rouges et 5 chaussettes bleues, mélangées. Dans le noir, combien faut-il en prendre au minimum pour être certain d’avoir deux chaussettes de la même couleur ? Pour être certain d’avoir deux chaussettes rouges ?\par\vspace{5mm}{\otLight\fontsize{8}{10}\selectfont\color{otMuted}Espace de recherche}\par\par\vspace{7mm}{\color{otRule!65}\hrule height .25pt}\par\vspace{7mm}{\color{otRule!65}\hrule height .25pt}\par\vspace{7mm}{\color{otRule!65}\hrule height .25pt}\par\vspace{7mm}{\color{otRule!65}\hrule height .25pt}\par\vspace{7mm}{\color{otRule!65}\hrule height .25pt}\par\vspace{7mm}{\color{otRule!65}\hrule height .25pt}
\end{blocEdition}
\editionPage{defi-ex-21}{D21 · Défi}{La spirale du robot}{Approfondissement · 20 à 45 minutes indicatives}
\begin{blocEdition}{defi-ex-21}
Un robot part du point O, tourné vers la droite. Il avance de 1 case, tourne à gauche, avance de 1 case, tourne à gauche, avance de 2 cases, tourne à gauche, avance de 2 cases, et ainsi de suite : deux fois chaque longueur, en augmentant de 1 à chaque paire. Il s’arrête après avoir parcouru la seconde ligne de 5 cases. Quelle distance a-t-il parcourue ? Où se trouve-t-il ?\label{fig-defi-21}\figureGeo{defi-spirale}
\par\vspace{5mm}{\otLight\fontsize{8}{10}\selectfont\color{otMuted}Espace de recherche}\par\par\vspace{7mm}{\color{otRule!65}\hrule height .25pt}\par\vspace{7mm}{\color{otRule!65}\hrule height .25pt}\par\vspace{7mm}{\color{otRule!65}\hrule height .25pt}\par\vspace{7mm}{\color{otRule!65}\hrule height .25pt}\par\vspace{7mm}{\color{otRule!65}\hrule height .25pt}\par\vspace{7mm}{\color{otRule!65}\hrule height .25pt}
\end{blocEdition}
